\ncuAbstractPage{\ncuCNtitle}{中文摘要}{摘\quad 要}
時間變化與類別不平衡常同時出現在分類問題中：歷史資料樣本較多，卻未必符合近期情況；近期資料較接近預測時點，少數類樣本卻較少。本研究以兩階段的集成更新模擬持續學習：保留歷史（Old）模型，整合近期（New）資料建立的新模型，並比較集成規則與特徵選取的影響。受資料可驗證範圍與本次配置限制，實證以固定 Old／New 分區評估一次模型池擴充。

實證使用 1999--2018 年美國企業的公開年度財務資料，以破產公司的最後觀測年度回溯重建下一年度事件代理標籤，未逐筆核對申請日；正類僅占全部紀錄的 0.77\%。實驗一在相同年份配置下比較四種分類器的 Old、New 與合併重訓（Retrain），另暫列 TabM 與 TabICL 的既有結果，待事件數據更新；框架依驗證表現選定最佳基礎模型，目前先以 XGBoost 建立 Old／New 六模型池測試流程，比較靜態集成、動態分類器選擇與動態集成選擇；實驗二在採樣前加入四種特徵選取方法與兩種保留數。所有成員選擇與分類門檻均由驗證資料決定。

目前測試顯示，近期資料的優勢依分類器而異。在目前測試模型池中，只整合近期模型的 New3 平均 AUC 最高（0.8601），依驗證資料選出兩個成員的靜態集成 F1-score 最高（0.1795）；全部六模型等權集成與動態方法未超過上述 AUC 與 F1-score 平均值。特徵選取沒有在多數方法與指標上一致改善。公司群聚配對分析顯示，New3 與驗證選定靜態集成相對六模型等權集成的 AUC 增加通過 165 項 Holm 校正，但後者的 F1-score 增加未通過校正。調參後的單一分類器亦有接近或高於集成的描述性分數；因參數與驗證資料不同，本研究尚不能主張集成優於單一模型。改用公司最終狀態標籤時，方法的排序隨之改變，顯示標籤定義會影響結論。

\noindent\textbf{關鍵字：}持續學習、類別不平衡、時間資料、模型集成、動態選擇、特徵選取、破產預測
